%! suppress = EscapeAmpersand
%! suppress = TooLargeSection
%! suppress = MissingLabel

\documentclass[10pt]{SPIIRAS_Proceedings}

\usepackage{mathtools}
\usepackage{proof}
\usepackage{booktabs}
\usepackage{array}
\let\ditshape\itshape

\newcommand{\vor}{~|~}
\newcommand{\step}{\longmapsto}
\newcommand{\hstep}{\mathrel{\rightsquigarrow}}
\newcommand{\keyword}[1]{\mathbf{#1}}

% --- notation for the calculus (Section 3) ---
\newcommand{\free}{\mathsf{free}}
\newcommand{\local}{\mathsf{local}}
\newcommand{\Any}{\mathsf{Any}}
\newcommand{\tj}[3]{#1 \vdash #2 : #3}
\newcommand{\lsub}[3]{#1 \vdash_{\ell} #2 <: #3}
\newcommand{\tsub}[3]{#1 \vdash #2 <: #3}
\newcommand{\cl}{\mathsf{cl}}
\newcommand{\ltof}{\ell}
\newcommand{\noloc}{\mathsf{noLocal}}
\newcommand{\elimp}{\mathsf{elim}^{+}}
\newcommand{\cautomaton}{}
\newcommand{\arr}[1]{\xrightarrow{\,#1\,}}
\renewcommand{\cap}{\keyword{cap}} % shadow TeX's \cap (set intersection, unused here)
\newcommand{\reic}{\keyword{resume}}
\newcommand{\hdlr}{\keyword{handler}}
\newcommand{\rname}[1]{\textsc{#1}}

\usepackage{listings}
\lstdefinestyle{modern}{
    basicstyle=\fontsize{9pt}{10pt}\ttfamily,
    keywordstyle=\bfseries,
    stringstyle=,
    commentstyle=\itshape,
    numberstyle=\tiny,
    tabsize=4,
    showspaces=false,
    showstringspaces=false,
    breaklines=true,
    breakatwhitespace=true,
    captionpos=b,
    escapeinside={(*@}{@*)},
    gobble=8,
}
\lstdefinelanguage{colang}{
    keywords=[1]{let,var,effect,perform,data,context,fun,where,if,else,throw,try,catch,local,free,op,handle,resume,return},
    sensitive=true,
    comment=[l]{//},
    morecomment=[s]{/*}{*/},
    morestring=[b]",
}
\lstset{style=modern}

% Helper macros for an English-primary layout that reuses the journal class
% without modifying SPIIRAS_Proceedings.cls. \udktop prints the UDC (+ optional
% DOI) line once at the very top; \makerustitle reproduces the Russian title
% block of \maketitle but WITHOUT the UDC line, so UDC is never duplicated.
\makeatletter
\newcommand{\udktop}{\noindent\footnotesize\textbf{\foreignlanguage{russian}{УДК}}~\@udk\ifx\@doi\@empty\relax\else\hfill\footnotesize\textbf{DOI}~\@doi\fi}
\newcommand{\makerustitle}{%
\begin{center}
\normalsize\@authorsTitleRus\\
\normalsize\bfseries\MakeUppercase{\@titleRus}\\[8pt]
\normalsize\normalfont
\end{center}
\rule{\textwidth}{0.5pt}
\textit{\@authorsRus}~\textbf{\@titleRus.}\\
{\hspace*{0.5cm}\textbf{Аннотация.}}~{\hyphenpenalty=10000\setlength{\parindent}{0.5cm} \@abstractRus}\\
\hspace*{0.5cm}\textbf{Ключевые слова:}~\@keywordsRus \\[-4pt]
\rule{\textwidth}{0.5pt} \\[-2pt]
}
\makeatother

\begin{document}
\selectlanguage{english}
% babel's Russian main language sets \lstlistingname to Cyrillic "Листинг" at
% \begin{document}, and \captionsenglish does not reset it; the paper's
% captions are English, so override here (after language selection).
\renewcommand{\lstlistingname}{Listing}

\udk{004.052.42, 004.43}
\titleRus{Машинно-проверенная метатеория исчисления escape-анализа с экзистенциальными временами жизни и обработчиками эффектов}
\authorsTitleRus{А. С\smallcapsfake{тоян}}
\authorsRus{Стоян~А.}
\titleEng{A mechanized metatheory of an escape-analysis calculus with existential lifetimes and effect handlers}
\authorsTitleEng{A. S\smallcapsfake{toyan}}
\authorsEng{Stoyan~A.}

\abstractRus{Типизированный escape-анализ переносит вопрос о выходе значения за пределы его лексической области видимости из оптимизирующего прохода компилятора в систему типов.
В работе рассматривается исчисление Core-Delta, в котором времена жизни образуют двухточечную решётку с операцией минимума (формально это join решётки), составные данные скрывают времена жизни полей экзистенциально и восстанавливают их при сопоставлении с образцом, а алгебраические обработчики эффектов передаются как значения-возможности.
Цель~--- дать машинно-проверенную метатеорию этого исчисления, точно установить, какие гарантии escape-анализа следуют из правил типизации, и тем самым дать компиляторам, совмещающим стековое или регионное размещение значений с обработчиками эффектов, верифицированную спецификацию типизации.
Приводится формальная система: синтаксис, подтипизация времён жизни и типов, типизация с правилом минимума для конструкторов и устранением экзистенциальных времён жизни с учётом вариантности, а также семантика редукций с runtime-маркерами, делимитерами и реифицированными продолжениями.
Синтаксис, семантика и доказательства механизированы в Rocq/Coq на чистых индексах де Брёйна (около 29 тыс. строк, без аксиом, незавершённых доказательств и гипотез).
Доказаны продвижение, сохранение типа и типобезопасность: все runtime-инварианты маркерной семантики машинно доказаны сохраняющимися при редукции.
Для исходных программ без каких-либо предположений о runtime-инвариантах доказаны три гарантии escape-анализа: каждое синтаксическое вхождение возможности имеет в области видимости делимитер со своим маркером, каждое выполненное событие на двух охраняемых каналах границы обработчика несёт значение нелокального типа, а значение, доставленное по нелокальному типу, не содержит возможностей и делимитеров ни на какой глубине.
Библиотека примеров содержит обработчики reader, state, exception и optionality, а также машинно-проверенные теоремы об отклонении целых программ, нарушающих локальность.
Ни одна ключевая теорема не зависит от недоказанных допущений: доверенная база сводится к ядру Rocq и его стандартной библиотеке.
Механизация переиспользуема как спецификация типизации; размещение в памяти остаётся отдельным аргументом.}

\keywordsRus{системы типов, escape-анализ, экзистенциальные времена жизни, алгебраические обработчики эффектов, ограничение области действия возможностей, машинно-проверенная метатеория, Rocq/Coq.}

\abstractEng{Type-based escape analysis moves the question of whether a value outlives its lexical scope from an optimizing compiler pass into the type system.
We study Core-Delta, a calculus in which lifetimes form a two-point lattice with a minimum operation (formally the lattice join), compound data hide field lifetimes existentially and recover them on pattern matching, and algebraic effect handlers are passed as capability values.
The goal is a machine-checked metatheory of this calculus that states precisely which escape-analysis guarantees follow from the typing rules, giving compilers that combine stack or region allocation with effect handlers a verified typing specification.
We give the formal system: syntax, lifetime and type subtyping, typing with a minimum rule for constructors and variance-aware elimination of existential lifetimes, and a reduction semantics with runtime markers, delimiters, and reified continuations.
Syntax, semantics, and proofs are mechanized in Rocq/Coq using pure de Bruijn indices (about 29 thousand lines, with no axioms, admitted proofs, or conjectures).
We prove progress, preservation, and type safety: every runtime invariant of the marker semantics is machine-checked to be preserved by reduction.
Three escape-analysis guarantees follow for source programs with no runtime-invariant assumptions: every capability occurrence has a delimiter carrying its marker in scope, every event on the two guarded boundary channels carries an escapably typed value, and values delivered at escapable types contain no capability or delimiter at any depth.
An example library exercises reader, state, exception, and optionality handlers, alongside machine-checked rejection theorems for complete leaking programs.
Every capstone theorem depends on no unproved assumptions: the trusted base is the Rocq kernel and its standard library.
The mechanization is reusable as a typing specification; memory placement remains a separate argument.}

\keywordsEng{type systems, escape analysis, existential lifetimes, algebraic effect handlers, capability confinement, mechanized metatheory, Rocq/Coq.}

\udktop
\makeengtitle
\normalsize

\section{Introduction}

Escape analysis determines whether a value can outlive its lexical scope.
It originated as a compiler optimization, for example in stack allocation and the removal of intermediate objects~\cite{park1992escape,blanchet1999escape}---and an optimization can afford to be whole-program and best-effort.
Lifting the same question into types puts escape information into function signatures, and both consequences follow from that one move: the analysis becomes \emph{modular}---each function is checked once, against its callees' signatures---and non-escape becomes part of the \emph{interface}, something a programmer can state and a language can promise, rather than an internal fact of the optimizer.
This is why modern languages move it into the type discipline: Rust ties it to lifetimes and borrowing~\cite{matsakis2014rust}, and OCaml develops a modal treatment of locality~\cite{lorenzen2024oxidizing}, recently extended to data-race freedom~\cite{georges2025data}; Kotlin illustrates the demand from outside the type discipline: a function can declare, as a \emph{contract}, that a lambda argument is invoked in place and does not escape the call---but the declaration is trusted, not verified~\cite{akhin2021kotlin}.

The original $Core_\Delta$ paper proposed a deliberately light discipline for this~\cite{stoyan2025existential}: lifetimes form a two-point lattice $\free \sqsubseteq \local$ with an explicit minimum, and---the key convenience---compound data types hide the lifetimes of their fields \emph{existentially}, recovering them at pattern matching, so user-defined types need not be parameterized by lifetimes.
Its design, subtyping, and a prototype type-checker are presented there.

A discipline meant to license stack or region placement, however, is only as good as its metatheory, and this metatheory is delicate: existential lifetimes interact with substitution, subtyping, and variance, and a pen-and-paper proof can easily miss a leak---a $\local$ value returned inside a closure's environment, an existential lifetime trapped in the result of a match.
The stakes set the proof standard: an unsound rule fails silently, corrupting memory only on the programs that exercise it.
The field's bar for such mechanisms is already machine-checked: RustBelt provides a semantic foundation for the safety of Rust~\cite{jung2018rustbelt}, complemented by the Stacked Borrows aliasing model~\cite{jung2020stacked}, while Prusti, Creusot, RustHornBelt, and RefinedRust lift properties of Rust programs to machine-checked specifications~\cite{astrauskas2022prusti,denis2022creusot,matsushita2022rusthornbelt,gaher2024refinedrust}.
Mechanization also returns more than confidence: it forces every rule into a precise, checked definition, and the artifact additionally makes the reduction semantics and the escape-specific side conditions executable---a certified evaluator and reflected deciders a compiler can be tested against---though the declarative typing relation itself is not algorithmic (Section~\ref{sec:discussion}).
The present work is the proof-side companion of the original paper: a machine-checked metatheory stating exactly which escape-analysis guarantees follow from the typing rules.

We prove the discipline not for the original calculus alone but extended with algebraic effect handlers, for three reasons.
Handlers are the \emph{sharpest instance} of scope-tied meaning: in a capability-passing discipline an operation is invoked through a capability value, valid only while its handler is on the stack, and first-class resumptions move values across scope boundaries in ways no local syntactic check can follow.
They are \emph{minimal and representative}: one construct covers exceptions, mutable state, generators, and asynchrony~\cite{plotkin2013handling}, and the discipline itself is not handler-specific---the same $\free$/$\local$ machinery governs closures, data, and any capability-like resource.
And they are where the field is active: Effekt, capture tracking, and evidence-passing handlers attack capability escape with explicit capture sets or a substantial polymorphic infrastructure~\cite{brachthauser2020effects,brachthauser2022effects,boruch2023capturing,xie2021generalized,xie2022first}; $Core_\Delta$'s coarser two-point discipline is the compact alternative, and the machine-checked guarantees are its warrant.
Section~\ref{sec:overview} tours the analysis informally before the formal development.

\paragraph{Model assumptions.}
Two assumptions keep the model concise and machine-checkable.
First, we work purely at the type level: there is no heap, store, or mutable reference, so ``escape'' is a lexical-lifetime discipline and stack allocation, ownership, and region placement are motivating applications rather than memory properties we prove.
Second, lifetimes form the two-point lattice $\free \sqsubseteq \local$; multi-level regions are a natural generalization.
Within these bounds the result is a faithful, executable, and checkable specification.

\paragraph{Contributions.}
The organizing principle is \emph{economy of mechanism}: handlers are the single effect construct, and they add no escape-analysis structure \emph{to the type language}---only $\noloc$ premises over the existing lattice, plus runtime machinery (markers and a scoping invariant).
First, we extend the handler-free $Core_\Delta$ with multi-operation, capability-passing effect handlers and give a reduction semantics with runtime markers, delimiters, and reified continuations (Section~\ref{sec:calculus}).
Second, we mechanize the entire system in Rocq/Coq with pure de Bruijn indices: syntax, lifetime and type subtyping, typing---including the constructor \emph{minimum} rule and the variance-aware elimination of existential lifetimes at \texttt{match}---and the reduction semantics.
Third, we prove progress, preservation, and type safety with \emph{no runtime-invariant assumptions for source programs}: every runtime invariant of the marker semantics is derived from source syntax and proved preserved by reduction (Section~\ref{sec:mechanization}).
Fourth, we establish three escape-analysis guarantees as machine-checked statements about the \emph{dynamics}, all holding for source programs with no runtime-invariant assumptions: occurrence-level capability confinement, per-event guarded-channel typing, and depth non-escape of runtime forms at escapable types; together with type safety they are bundled into one record and one theorem, displayed in Section~\ref{sec:mechanization}.
Fifth, we provide an executable example library and machine-checked \emph{negative} witnesses that reject locality-violating programs (Section~\ref{sec:evaluation}).
None of this metatheory---progress, preservation, or any escape guarantee---appears in the original paper.

\section{\texorpdfstring{$Core_\Delta$}{Core-Delta} by Example}\label{sec:overview}

Before the formal development, we sketch the analysis informally, in the surface notation of the original paper~\cite{stoyan2025existential}.
Types carry \emph{lifetime} labels: a type \lstinline[language=colang]|T| with lifetime \lstinline[language=colang]|lab| is written \lstinline[language=colang]{T'lab}; the label \lstinline[language=colang]|free| marks values that may escape, \lstinline[language=colang]|local| marks values confined to their defining scope, and a function type carries the lifetime of its \emph{closure}; a type whose overall lifetime restriction is \lstinline[language=colang]|free| we call \emph{escapable}.
A compound value's lifetime is the minimum (written \lstinline[language=colang]|+|) of its components' lifetimes; that paper develops this value discipline in detail, together with lifetime polymorphism and bounded quantification over tracked types.
This paper extends the same discipline to algebraic effect handlers.

The design's distinctive move---the original paper's namesake---is its \emph{existential} treatment of those component lifetimes.
User-defined types take no lifetime parameters; a constructor instead absorbs the parts' lifetimes into the annotation of the whole:
\begin{lstlisting}[language=colang]
        fun makeRepo(file: File'lf, conn: Connection'lc): Repository'lf+lc =
            Repository(file, conn)
\end{lstlisting}
The resulting type never says which part contributed which lifetime---the parts' individual lifetimes vanish from it, which is what makes them existential---yet nothing is lost for safety: a repository built around a $\local$ connection is itself $\local$, so packing a tracked value into a container cannot launder it.
Unpacking is the mirror image: matching a \lstinline[language=colang]|Repository'local| open recovers the connection's lifetime only as a fresh variable bounded by the repository's own $\local$, and a variance-aware approximation then eliminates that variable from the branch's result type (\rname{T-Match}, Section~\ref{sec:calculus}).

Entering a $\keyword{handle}$ block mints a \emph{capability}---a first-class value granting the right to perform the effect's operations---and passes it to the handled code; a capability is meaningful only while its handler is on the stack, so an escaped capability would address a handler that no longer exists.
Preventing that escape needs no new machinery---a capability is simply a value whose data type sits at lifetime $\local$---and Listing~\ref{lst:state} shows the whole discipline at work on a two-operation state effect, written in state-passing style.

\begin{lstlisting}[language=colang,label={lst:state},caption={A two-operation state effect in state-passing style. One capability \texttt{st} grants both operations; the rejected \texttt{leak} tries to return it.}]
        effect State<s> { get: (Unit) -> s, put: (s) -> Unit }

        fun withState<s <: Any'free, r <: Any'free>(
                body: (State<s>'local) -local-> r): (s) -local-> r =
            handle st: State<s>'local {
                get()  => fun(s) => resume(s)(s)
                put(x) => fun(s) => resume(())(x)
            } in
                let res = body(st) in fun(s) => res

        withState(fun(st) =>
            let a = perform st.get() in    // reads 2
            let _ = perform st.put(3) in
            let b = perform st.get() in    // reads 3
            a + b)(2)                      // evaluates to 5

        // rejected: the body would answer at State<Int>'local,
        // and the handle rule demands an escapable answer type
        fun leak(): State<Int>'local =
            handle st: State<Int>'local { ... } in st
\end{lstlisting}

The effect declares two operations; $\keyword{handle}$ installs one clause per declared operation and hands the fresh capability \texttt{st} to \texttt{body}; a $\keyword{perform}$ names the capability and the operation it invokes.
The clauses use the classic state-passing encoding of mutable state~\cite{plotkin2013handling}: the $\keyword{handle}$'s answer is a function of the current state, and each \lstinline[language=colang]|resume| threads it onward.
The closing call shows the observable behaviour: started at \lstinline[language=colang]|2|, the run answers the first \lstinline[language=colang]|get()| with \lstinline[language=colang]|2|, makes \lstinline[language=colang]|put(3)| visible to the second \lstinline[language=colang]|get()|, and delivers the sum \lstinline[language=colang]|5|.
A clause that never resumes recovers exceptions, one that resumes twice expresses backtracking, and nothing in the analysis depends on how many operations an effect declares.

The escape analysis itself needs nothing new for any of this.
Safety demands one thing: nothing that can still reach the capability may leave the $\keyword{handle}$ \emph{escapably typed}.
The discipline delivers it from two sides.
On the value side, locality is contagious: \texttt{st} is an ordinary $\local$ value, so by the minimum rule a closure capturing it or a container storing it is $\local$ in turn, at any depth---which is also why \texttt{body} must declare its argument $\local$.
On the boundary side, everything the type system lets cross out---the body's answer, checked by the $\keyword{handle}$ rule, and likewise the arguments operations carry to their clauses---must be escapable.
Together they close every escapable exit: the final \texttt{leak} would answer at \lstinline[language=colang]{State<Int>'local} and is rejected, and wrapping \texttt{st} in an \lstinline[language=colang]|Option| changes nothing---that container check is exactly our first negative witness (Section~\ref{sec:evaluation}).
Flows at \emph{local} types may still cross---a handler may even hand its resumption to the enclosing scope---but such values stay confined at non-escapable types, where the same discipline keeps governing them (Section~\ref{sec:mechanization}).

The guarantees we prove can be read informally.
First, a capability---every one of its syntactic occurrences, even under $\lambda$s and inside data---always sits under its handler.
Second, the two data channels crossing the handler boundary (the operation argument going in, the delimited body's answer coming out) only ever carry values of escapable types.
Third, a program typed at an escapable type only ever delivers values containing no capability or \emph{running} handler at any depth.
Section~\ref{sec:mechanization} states all three as machine-checked theorems over the reduction relation.

\section{Related Work and Problem Statement}\label{sec:related-work}

Classical region systems made lifetime an explicit object of typing long before modern ownership languages~\cite{tofte1997region}.
Their aim is close to ours: to replace a global analysis with local conditions at function boundaries.
Type-based \emph{escape} analysis itself goes back to Hannan~\cite{hannan1998type}, who reformulates escape inference for functional languages as a type system.
The difference is that $Core_\Delta$ does not parameterize user types by lifetimes: a container's restriction is derived from its fields and recovered existentially at match (Section~\ref{sec:overview}), which trades annotation burden for a harder proof of preservation under reduction.

Effect handlers have been studied denotationally and operationally~\cite{plotkin2013handling,zyuzin2021contextual,van2022handling,van2024framework,tang2025modal}, integrated into OCaml, WebAssembly, and Unison's abilities~\cite{sivaramakrishnan2021retrofitting,phipps2023continuing,unison-abilities}, and compiled via evidence passing~\cite{xie2021generalized}.
Where linear and affine disciplines police \emph{how} continuations are used~\cite{tang2024soundly,vanrooij2025affect}, $Core_\Delta$ polices \emph{which} values may cross the handler boundary.
Closest are the capability-based systems---Effekt's scope- and type-based reasoning~\cite{brachthauser2020effects,brachthauser2022effects} and capture tracking~\cite{boruch2023capturing}; a modal framework compares row- and capability-based systems in one setting~\cite{tang2026modal}.
$Core_\Delta$ records coarser information---a value is either free to escape or locally confined---and, unlike systems that extend the type language with rows or capture sets, adds no effect-specific type structure.

Capability confinement itself has many routes: lexically scoped handler instances~\cite{biernacki2019binders} and scoped and higher-order operations~\cite{bachpoulsen2023hefty}; scoped capabilities in Scala~\cite{odersky2022scoped}; second-class values~\cite{osvald2016gentrification}; effect rows~\cite{leijen2017type}; evidence- and tunnel-based designs~\cite{zhang2020handling}; and mechanized capture calculi with explicit boxes~\cite{xu2023formalizing} or without them~\cite{xu2025s}---a broader survey situates these lines~\cite{stoyan2024modern}.
Mainstream practice, by contrast, hands capabilities around through the typing context---Scala's safer exceptions~\cite{odersky2021safer}, context parameters in Kotlin~\cite{KEEP-context-parameters}---which controls who may \emph{obtain} a capability but not whether it \emph{outlives} its scope; escape analysis supplies exactly the missing half.
$Core_\Delta$ folds confinement into the two-point value-lifetime discipline it already uses for ordinary data.

Mechanizing type-system metatheory builds on POPLMark and the engineering of formal metatheory~\cite{aydemir2005poplmark,aydemir2008engineering}, with syntax and substitution automation available~\cite{stark2019autosubst2}; handler metatheory exists in separation-logic form~\cite{devilhena2021separation}.
We deliberately use pure de Bruijn indices with ordinary recursive shift and substitution functions: transparent and portable, at the cost of an explicit substitution metatheory.
The alternatives fit poorly: $Core_\Delta$'s constructors, $\keyword{match}$ branches, and handler clauses bind whole lists of lifetimes, types, and terms at once---variable-arity binding that sits outside the comfortable fragment of the automation above.

$Core_\Delta$ targets the combination these lines leave open, and the model must therefore (i) express lifetimes in types without burdening the user with lifetime parameters; (ii) confine a runtime capability to its handler; (iii) admit first-class continuations while keeping them out of every escapably typed position; and (iv) have a checkable metatheory suitable as a compiler specification (Sections~\ref{sec:mechanization}--\ref{sec:evaluation}).

\section{The \texorpdfstring{$Core_\Delta$}{Core-Delta} Calculus}\label{sec:calculus}

The mechanization covers the system in full; the figures of this section present the core rules, with routine well-formedness, length, and effect-lookup premises elided.

\subsection{Syntax}

Figure~\ref{fig:syntax} gives the grammar.
Lifetimes are the two-point lattice: $\Delta_1+\Delta_2$ denotes the lifetime \emph{minimum}, which is $\local$ as soon as either operand is $\local$.
In the order $\free \sqsubseteq \local$ the more restrictive (shorter) lifetime sits \emph{higher}, so this minimum is formally the least upper bound.
A function type $\tau \arr{\Delta} \sigma$ carries a \emph{closure lifetime} $\Delta$; a data type $K\,\Delta\,\overline{\tau}$ is annotated with the lifetime for which a $K$-value may be used.
Effect capabilities reuse the data machinery: a capability for effect $E$ has type $E\,\local\,\overline{\tau}$.
Core functions are unary; multi-argument functions and surface effect declarations are sugar; type and lifetime abstractions are restricted so that every $\Lambda$-chain ends in a $\lambda$.

Two runtime-only constructors extend the grammar during reduction and close it under the semantics: the capability value $\cap_E^{m,\sigma}$ and the delimiter $\hdlr_m$, each carrying a runtime marker $m$; a \emph{source} term is one containing neither.
Type annotations ($\sigma$ is the handler's public answer type) are recorded so that preservation can re-type the reducts, and the mechanized source $\keyword{handle}$ and $\keyword{perform}$ forms carry answer- and result-type annotations (elided here) for the same reason.
The reified resumption $\reic_m^{\,\tau,\sigma}(x.\,t)$ is \emph{derived notation}, not a constructor: it abbreviates the delimiter-re-installing abstraction $\lambda(x{:}A).\,\hdlr_m^{\,\tau,\sigma}\,t$ (Figure~\ref{fig:handlers-sem}, caption), and is therefore a value because $\lambda$-abstractions are.

\begin{figure}[h]
\footnotesize
\setlength{\arraycolsep}{2pt}
\[
\begin{array}{r@{\;}c@{\;}l}
\Delta &::=& l \vor \free \vor \local \vor \Delta_1 + \Delta_2 \\[2pt]
\tau,\sigma &::=& \alpha \vor \tau \arr{\Delta} \sigma \vor K\,\Delta\,\overline{\tau}
   \vor \forall l.\,\tau \vor \forall(\alpha{<:}\tau).\,\tau \\[2pt]
t &::=& x \vor \lambda(x{:}\tau).\,t \vor t\;t
   \vor \Lambda(\alpha{<:}\tau).\,t \vor t\,[\tau]
   \vor \Lambda l.\,t \vor t\,[\Delta] \\
  && \vor\; K[\Delta;\overline{\Delta};\overline{\tau}](\overline{t})
   \vor \keyword{match}\;t\;\{K(\overline{x})\Rightarrow t \mid \_ \Rightarrow t\} \\
  && \vor\; \keyword{handle}\;E\,[\overline{\tau}]\,\{\overline{op\Rightarrow t}\}\;\keyword{in}\;t
   \vor \keyword{perform}\;t.op\,[\overline{\tau}]\,t \\
  && \vor\; \cap_E^{m,\sigma}[\overline{\tau}]\{\overline{op\Rightarrow t}\}
   \vor \hdlr_m^{\,\tau,\sigma}\,t \\[2pt]
v &::=& \lambda(x{:}\tau).t \vor \Lambda(\alpha{<:}\tau).t \vor \Lambda l.t
   \vor K[\dots](\overline{v}) \vor \cap_E^{m,\sigma}[\dots]
\end{array}
\]
\caption{Syntax of $Core_\Delta$ with handlers (de Bruijn in the artifact). One grammar covers source and runtime terms: source terms contain no $\cap$/$\hdlr$, and $\reic$ is derived notation for a $\lambda$-abstraction.}
\label{fig:syntax}
\end{figure}

\subsection{Lifetimes and subtyping}

Lifetime subtyping (Figure~\ref{fig:sub}, top) makes $\free$ the bottom and $\local$ the top, with $\Delta_1+\Delta_2$ the least upper bound: together with the omitted reflexivity and transitivity, the rules place it above both operands (\rname{LS-JoinR}) and below any common upper bound (\rname{LS-JoinL}).
The lifetime restriction $\ltof$ is computed structurally: $\ltof(\alpha)=\free$ (the $\Gamma$-aware $\ltof_\Gamma$ consults the variable's bound), $\ltof(\tau\arr{\Delta}\sigma)=\Delta$---the closure lifetime only, the domain and codomain are \emph{not} inspected---$\ltof(K\,\Delta\,\overline{\tau})=\Delta+\ltof(\overline{\tau})$, and $\ltof(\forall\ldots)=\local$.
The abbreviation $\noloc_\Gamma(\tau)$ for $\lsub{\Gamma}{\ltof_\Gamma(\tau)}{\free}$ is the formal reading of ``escapable'' from Section~\ref{sec:overview} and recurs throughout the typing rules and theorems.
Type subtyping (bottom) is covariant in a data type's lifetime and invariant in its type arguments (\rname{S-Data}); contravariant in a function's domain and covariant in its closure lifetime and codomain (\rname{S-Fun}); it also includes an $\Any$ rule by which any $\tau$ upcasts to $\Any\,\Delta$---in the artifact a data type with a distinguished tag---provided $\ltof_\Gamma(\tau) <: \Delta$.
$\Any$ has no elimination form (\rname{T-Match} excludes its tag), so it is a top-like sink: it states $\ltof$ bounds, and $\Any\,\local$ serves as a deliberately imprecise interface in the precision examples of Section~\ref{sec:evaluation}.
Bounded quantifiers follow full $F_{<:}$: the bound is contravariant, and the bodies are compared under the tighter bound.

\begin{figure}[h]
\footnotesize
\centering
$\infer[\rname{LS-Free}]{\lsub{\Gamma}{\free}{\Delta}}{}$
\quad
$\infer[\rname{LS-Local}]{\lsub{\Gamma}{\Delta}{\local}}{}$
\quad
$\infer[\rname{LS-JoinL}]{\lsub{\Gamma}{\Delta_1+\Delta_2}{\Delta}}{\lsub{\Gamma}{\Delta_1}{\Delta} & \lsub{\Gamma}{\Delta_2}{\Delta}}$
\quad
$\infer[\rname{LS-JoinR}]{\lsub{\Gamma}{\Delta}{\Delta_1+\Delta_2}}{\lsub{\Gamma}{\Delta}{\Delta_i}}$
\\[1.6ex]
$\infer[\rname{S-Data}]{\tsub{\Gamma}{K\,\Delta\,\overline{\tau}}{K\,\Delta'\,\overline{\tau}}}{\lsub{\Gamma}{\Delta}{\Delta'}}$
\quad
$\infer[\rname{S-Fun}]{\tsub{\Gamma}{A'\arr{\Delta}B}{A\arr{\Delta'}B'}}{\tsub{\Gamma}{A}{A'} & \lsub{\Gamma}{\Delta}{\Delta'} & \tsub{\Gamma}{B}{B'}}$
\quad
$\infer[\rname{S-Any}]{\tsub{\Gamma}{\tau}{\Any\,\Delta}}{\lsub{\Gamma}{\ltof_\Gamma(\tau)}{\Delta}}$
\caption{Lifetime subtyping (top) and selected type subtyping (bottom); reflexivity, transitivity, variable, and quantifier rules omitted.}
\label{fig:sub}
\end{figure}

\subsection{Typing}

Figure~\ref{fig:typing} shows the core typing rules.
The escape discipline is concentrated in three places.

\emph{Functions (\rname{T-Lam}).} A $\lambda$ receives a closure lifetime $\Delta$ that must bound $\cl_\Gamma(t)$, the \emph{capture lifetime} of its body.
$\cl_\Gamma(t)$ is $\local$ if the body syntactically contains a marker-bearing runtime form (a capability or a delimiter); otherwise it is the minimum ($+$) of the lifetimes $\ltof_\Gamma$ of the captured free variables' types.
Thus a function that captures a $\local$ value (or a capability) necessarily has closure lifetime $\local$, and by \rname{S-Fun} and $\ltof$ it can no longer be embedded where a $\free$ lifetime is required.
This is how a local value is kept from crossing a function boundary, without any separate ``result must be non-local'' side condition.

\emph{Constructors (\rname{T-Ctor}).} A constructor's result lifetime must dominate the minimum of its field lifetimes: $\lsub{\Gamma}{\ltof(\overline{\rho})}{\ltof(K\,\Delta\,\overline{\tau})}$, where $\overline{\rho}$ are the instantiated field types.
A value built from fields is therefore safe for no longer than its shortest-lived field.
A disjointness side condition keeps data tags apart from effect tags; the $\Any$ tag is excluded at elimination (\rname{T-Match}).

\emph{Pattern matching (\rname{T-Match}).} Matching a scrutinee of type $K\,\Delta\,\overline{\tau}$ introduces $n$ fresh lifetime variables $\overline{l}$ bounded by $\Delta$, types the success branch under the recovered field types to a type $\eta$, and then \emph{eliminates} the fresh variables from $\eta$ with $\elimp_{\overline{l}\mapsto\Delta}$: a variable in a positive position is replaced by its bound $\Delta$, in a negative position by $\free$, and an occurrence in an invariant position is \emph{rejected}---$\elimp$ is a partial function, and the premise $\elimp_{\overline{l}\mapsto\Delta}(\eta)=\eta'$ of \rname{T-Match} requires it to succeed, so a match whose result type would trap an existential lifetime in an invariant position is ill-typed.
This recovers the existentially hidden lifetimes soundly and is exactly the rule that blocks leakage through containers and functions.

\begin{figure}[h]
\footnotesize
\centering
$\infer[\rname{T-Var}]{\tj{\Gamma}{x}{\tau}}{(x{:}\tau)\in\Gamma}$
\quad
$\infer[\rname{T-Sub}]{\tj{\Gamma}{t}{U}}{\tj{\Gamma}{t}{T} & \tsub{\Gamma}{T}{U}}$
\quad
$\infer[\rname{T-App}]{\tj{\Gamma}{t_1\,t_2}{B}}{\tj{\Gamma}{t_1}{A\arr{\Delta}B} & \tj{\Gamma}{t_2}{A}}$
\\[1.6ex]
$\infer[\rname{T-Lam}]{\tj{\Gamma}{\lambda(x{:}A).\,t}{A\arr{\Delta}B}}
  {\tj{\Gamma,x{:}A}{t}{B} & \lsub{\Gamma}{\cl_\Gamma(t)}{\Delta}}$
\qquad
$\infer[\rname{T-TyApp}]{\tj{\Gamma}{t\,[S]}{[\alpha{\mapsto}S]U}}
  {\tj{\Gamma}{t}{\forall(\alpha{<:}B).U} & \tsub{\Gamma}{S}{B}}$
\\[1.6ex]
$\infer[\rname{T-Ctor}]{\tj{\Gamma}{K[\Delta;\overline{\Delta};\overline{\tau}](\overline{t})}{K\,\Delta\,\overline{\tau}}}
  {\lsub{\Gamma}{\ltof(\overline{\rho})}{\ltof(K\,\Delta\,\overline{\tau})} & \overline{\lsub{\Gamma}{\Delta_i}{\Delta}} & \overline{\tj{\Gamma}{t_i}{\rho_i}}}$
\\[1.6ex]
$\infer[\rname{T-Match}]{\tj{\Gamma}{\keyword{match}\;t\;\{K(\overline{x})\Rightarrow t_y \mid \_\Rightarrow t_n\}}{\eta'}}
  {\tj{\Gamma}{t}{K\,\Delta\,\overline{\tau}} & \tj{\Gamma,\overline{l}{<:}\Delta,\overline{x{:}\rho}}{t_y}{\eta} & \elimp_{\overline{l}\mapsto\Delta}(\eta)=\eta' & \tj{\Gamma}{t_n}{\eta'}}$
\caption{Core typing rules (well-formedness and length premises elided).
$\cl$ is the capture lifetime; $\elimp$ is variance-aware existential-lifetime elimination.
A constructor's tag and its declared result-type tag may differ (as in \texttt{Some}/\texttt{Option}); the figure identifies them for readability.}
\label{fig:typing}
\end{figure}

\subsection{Handlers: typing and semantics}

A capability value has a $\local$ data type, and a resumption is a function with closure lifetime $\local$ (Figure~\ref{fig:handlers-typing}): both are therefore barred from every escapably typed position by the lifetime discipline above.
An effect declares a \emph{list} of operations, and a $\keyword{handle}$ supplies one clause per declared operation.
The source rule \rname{T-Handle} types \emph{each} clause $t_{op}$ against the public answer type $T_R$, with that operation's argument $x{:}\sigma^{op}_\beta$ and resumption $k{:}\rho^{op}_\beta\arr{\local}T_R$ in scope, requires the clause list to cover exactly the declared operations, and types the handled body $b$ against $T_B$ with the capability $c$ in scope; $\sigma^{op}_\beta,\rho^{op}_\beta$ are operation $op$'s declared argument/result types instantiated at the effect's type arguments $\overline{\tau}$; each operation's type parameters $\overline{\beta}$ are bound by its signature in the effect declaration, scope over $\sigma^{op}_\beta$, $\rho^{op}_\beta$, and its clause, and are instantiated by \rname{H-Perform} at the perform site.
The runtime forms mirror this: a capability records $T_R$ and types its stored operation clauses the same way (\rname{T-Cap}), and a handler delimiter is transparent to typing, requiring only that the body answer type $T_B$ be escapable and a subtype of the public answer $T_R$ (\rname{T-HandlerM}).
The \rname{T-Perform} rule selects an operation $op$ of the capability's effect and instantiates its signature once more with the perform-site type arguments $\overline{S}$, giving argument type $\sigma^{op}_{\overline{S}}$ and result type $\rho^{op}_{\overline{S}}$; it requires the argument type and each $S_i$ to be escapable ($\noloc_\Gamma$), so a value crossing the operation boundary cannot smuggle a local capability out of scope.
The operation \emph{result}, however, may be local, which lets a deep resumption deliver local values back to the operation site.

The reduction rules (Figure~\ref{fig:handlers-sem}) are layered: head reductions $t \hstep t'$ have no congruence, and $t \step t'$ closes them under well-formed call-by-value evaluation contexts $\mathcal{E}$ via $\rname{S-Step}$.
Entering a handler ($\rname{S-Handle}$) picks a fresh marker $m$, builds the capability $\cap_E^{m}$, substitutes it into the body, and installs the delimiter $\hdlr_m$.
A $\keyword{perform}$ of operation $op$ on a capability with marker $m$, sitting inside a context $P$ that is \emph{pure for $m$}---it contains no delimiter carrying $m$, delimiters for other markers being allowed---and a well-formed call-by-value evaluation context, under the matching $\hdlr_m$, captures $P$ as a resumption $\reic_m$ that re-installs the delimiter, and runs \emph{that operation's} stored clause on the argument and the resumption ($\rname{H-Perform}$).
Purity selects the matching delimiter and well-formedness forces left-to-right call-by-value capture; together they make the captured continuation unique, so one-step reduction is deterministic up to the globally fresh marker minted by \rname{S-Handle} (Section~\ref{sec:mechanization}).
Applying a resumption re-installs the delimiter around the resumed body ($\rname{H-Resume}$); the handlers are thus \emph{deep}---and since $\reic_m$ abbreviates exactly such a re-installing $\lambda$, shallow handlers are not expressible in this presentation.
The $\overline{\beta}$ binders of the selected operation carry the bound $\Any\,\free$---the static counterpart of \rname{T-Perform}'s $\overline{\noloc_\Gamma(S)}$---and in the mechanization the captured context is shifted across the resumption's new $\lambda$-binder.
The invariant $\mathsf{well\_scoped}\;\overline{m}\;t$ records that every runtime capability with marker $m$ lies under its $\hdlr_m$ or within the active set $\overline{m}$, together with a provenance condition on stored operation bodies; source programs satisfy $\mathsf{well\_scoped}\;[\,]\;t$.
The provenance form is forced: the two more obvious invariants---direct marker confinement and marker freshness---are not preserved by reduction (Section~\ref{sec:mechanization}).

\textit{Economy of mechanism.} Handlers add no escape-analysis type structure.
The type grammar of Figure~\ref{fig:syntax} is closed under the same formers with or without effects---there is no effect row, capture set, or effect-lifetime parameter.
An effect is registered as a context signature and denoted by a reserved tag on the data former $K\,\Delta\,\overline{\tau}$, kept disjoint from data tags by the side condition $\mathsf{lookupEff}(\Gamma,K)=\bot$ on \rname{T-Ctor}/\rname{T-Match}; a capability is then a data value at lifetime $\local$ and a resumption a $\local$ function.
Confinement and value non-escape are therefore governed by one lattice $\free\sqsubseteq\local$ and one elimination discipline, and the reader, state, exception, and optionality handlers of Section~\ref{sec:evaluation} are four instances of this one construct.
What handlers \emph{do} add is operational---runtime markers, the delimiter $\hdlr_m$, and the $\mathsf{well\_scoped}$ invariant; the economy is at the level of types and escape analysis, not the runtime.

\begin{figure}[!htp]
\footnotesize
\centering
$\infer[\rname{T-Handle}]{\tj{\Gamma}{\keyword{handle}\;E\,[\overline{\tau}]\,\{\overline{op\Rightarrow t_{op}}\}\;\keyword{in}\;b}{T_R}}
  {\begin{array}{c}
   \overline{\tj{\Gamma,\overline{\beta{<:}\Any\,\free},x{:}\sigma^{op}_\beta,\,k{:}\rho^{op}_\beta\arr{\local}T_R}{t_{op}}{T_R}} \\[1pt]
   \tj{\Gamma,c{:}E\,\local\,\overline{\tau}}{b}{T_B} \quad
   \noloc_\Gamma(T_B) \quad
   \tsub{\Gamma}{T_B}{T_R}
   \end{array}}$
\\[1.6ex]
$\infer[\rname{T-Cap}]{\tj{\Gamma}{\cap_E^{m,T_R}[\overline{\tau}]\{\overline{op\Rightarrow t_{op}}\}}{E\,\local\,\overline{\tau}}}{\overline{\tj{\Gamma,\overline{\beta{<:}\Any\,\free},x{:}\sigma^{op}_\beta,\,k{:}\rho^{op}_\beta\arr{\local}T_R}{t_{op}}{T_R}}}$
\quad
$\infer[\rname{T-Resume}]{\tj{\Gamma}{\reic_m^{\,T_B,T_R}(x.t)}{A\arr{\local}T_R}}{\tj{\Gamma,x{:}A}{t}{T_B} & \noloc_\Gamma(T_B) & \tsub{\Gamma}{T_B}{T_R}}$
\\[1.6ex]
$\infer[\rname{T-Perform}]{\tj{\Gamma}{\keyword{perform}\;t.op\,[\overline{S}]\,a}{\rho^{op}_{\overline{S}}}}{\tj{\Gamma}{t}{E\,\Delta\,\overline{\tau}} & op\in E & \tj{\Gamma}{a}{\sigma^{op}_{\overline{S}}} & \noloc_\Gamma(\sigma^{op}_{\overline{S}}) & \overline{\noloc_\Gamma(S)}}$
\quad
$\infer[\rname{T-HandlerM}]{\tj{\Gamma}{\hdlr_m^{\,T_B,T_R}\,t}{T_R}}{\tj{\Gamma}{t}{T_B} & \noloc_\Gamma(T_B) & \tsub{\Gamma}{T_B}{T_R}}$
\caption{Handler typing.
Overlined premises and clauses range over the effect's declared operations, which the clause list must cover exactly; $op\in E$ abbreviates looking up $op$'s signature in $E$'s declaration; each clause is typed under its operation's type binders $\overline{\beta}$.
$\sigma^{op}_\beta,\rho^{op}_\beta$ (resp.\ $\sigma^{op}_{\overline{S}},\rho^{op}_{\overline{S}}$) are operation $op$'s argument/result types instantiated at the effect's type arguments $\overline{\tau}$ (resp.\ additionally at the perform-site arguments $\overline{S}$); in the artifact $T_R$ is additionally shifted across the $\overline{\beta}$ binders; effect lookup, well-formedness, and length premises are elided.}
\label{fig:handlers-typing}
\end{figure}

\begin{figure}[!htp]
\footnotesize
\centering
$(\lambda(x{:}T).t)\,v \hstep [x{\mapsto}v]t$
\qquad
$\hdlr_m^{\,T_B,T_R}\,v \hstep v$
\qquad
$\keyword{match}\;K[\Delta;\overline{\Delta'};\overline{\tau}](\overline{v})\;\{K(\overline{x})\Rightarrow t_y\mid\_\}\hstep [\overline{l}{\mapsto}\overline{\Delta'},\,\overline{x}{\mapsto}\overline{v}]\,t_y$
\\[1.2ex]
$\hdlr_m^{\,T_B,T_R}\big(P[\,\keyword{perform}\;\cap_E^{m}[\dots].op\,[\overline{S}]\,v\,]\big)\hstep [x{\mapsto}v,\,k{\mapsto}\reic_m^{\,T_B,T_R}(y.\,P[y])]\;[\overline{\beta}{\mapsto}\overline{S}]\;t_{op}$
\quad (\rname{H-Perform}; $P$ pure for $m$, $\mathsf{ectx\_wf}(P)$)
\\[1.2ex]
$(\reic_m^{\,T_B,T_R}(x.b))\,v \hstep \hdlr_m^{\,T_B,T_R}\,[x{\mapsto}v]b$
\quad (\rname{H-Resume})
\qquad
$\infer[\rname{S-Step}]{\mathcal{E}[q]\step \mathcal{E}[q']}{q\hstep q' & \mathsf{ectx\_wf}(\mathcal{E})}$
\qquad
$\infer[\rname{S-Handle}]{\mathcal{E}[\keyword{handle}\;E\,[\overline{\tau}]\,\{\overline{op{\Rightarrow}t_{op}}\}\;\keyword{in}\;b]\step \mathcal{E}[\hdlr_m^{\,T_B,T_R}\,([c{\mapsto}\cap_E^{m,T_R}]\,b)]}{m\ \text{globally fresh}}$
\caption{Reduction semantics of handlers.
In the match contraction, $\overline{l}$ are the branch's fresh lifetime variables of \rname{T-Match}, instantiated with the scrutinee's stored lifetimes $\overline{\Delta'}$.
In \rname{H-Perform}, $t_{op}$ is the stored clause that the index $op$ selects; purity picks the delimiter, well-formedness the unique call-by-value capture point; $\overline{\beta}$ are the selected operation's type parameters, instantiated at $\overline{S}$ before the term substitution.
In \rname{S-Handle} the minted marker is globally fresh (it occurs nowhere in the whole term), $c$ is the variable that $\keyword{handle}$ binds to the fresh capability in $b$, and the minted capability stores the clause list.
$\reic_m^{\,\tau,\sigma}(x.t)$ abbreviates $\lambda(x{:}A).\,\hdlr_m^{\,\tau,\sigma}\,t$, where $A$ is the resumption's argument type, so \rname{T-Resume} and \rname{H-Resume} are derived rules.}
\label{fig:handlers-sem}
\end{figure}

\section{Mechanization and Main Results}\label{sec:mechanization}

The development is in Rocq/Coq, using pure de Bruijn indices and ordinary recursive shift/substitution functions, and is published at \url{https://github.com/penguin-boxx/effect-proofs/tree/submission-iia-2026}.
The verified core is the modules of Table~\ref{tab:modules}, which build in dependency order---\texttt{core/} fixes the calculus, \texttt{meta/} proves the substitution metatheory, and \texttt{safety/} derives the guarantees; example files are separated and lie outside the dependency cone of the metatheory.

\begin{table}[h]
\centering
\caption{Structure of the Rocq development}
\label{tab:modules}
\footnotesize
\begin{tabular}{@{}p{0.34\linewidth}p{0.56\linewidth}@{}}
\toprule
\textbf{Package} & \textbf{Contents} \\
\midrule
\texttt{core/} & syntax, values, shifts and substitutions, evaluation contexts and the reduction semantics, markers, lifetime/type subtyping, capture lifetime, variance elimination, typing \\
\texttt{meta/} & shift and substitution commutation laws, weakening, narrowing, substitution-preserves-typing lemmas, closedness bookkeeping, variance, the evaluation-environment invariant (\texttt{eval\_ctx}) \\
\texttt{safety/} & marker invariants, canonical forms, progress, preservation, type safety, confinement (incl.\ occurrence paths), guarded-channel events, non-escape, certified deciders, marker renaming, determinism, executable semantics \\
\texttt{examples/} & surface notation, example handlers, reduction traces, machine-checked negative witnesses, complete rejection proofs with positive companions \\
\bottomrule
\end{tabular}
\end{table}

\subsection{Safety theorems}

We write $\step$ for one reduction step; theorems are stated over a typing environment $\Gamma$ containing only data-constructor and effect declarations---no term, type, or lifetime variables (\texttt{eval\_ctx} in the artifact)---which covers the empty environment of a closed program.
The soundness proof tracks runtime invariants alongside typing, bundled with the typing judgment as \texttt{safety\_invariants} in the artifact: $\mathsf{marker\_types\_safe}\;t$ (all delimiter and capability answer-type annotations for one marker agree), $\mathsf{marker\_annots\_closed}\;t$ (a technical strengthening that keeps those annotations stable under substitution), $\mathsf{well\_scoped}\;[\,]\;t$ (capabilities stay under their delimiters, strengthened with marker \emph{provenance}: each capability's stored operation body is well-scoped at the scope outside that capability's own delimiter), and $\mathsf{rt\_closed}\;t$ (stored operation bodies are closed apart from their own argument and resumption binders).
All hold trivially for source programs, which carry no runtime markers, and all are proved preserved by reduction; progress needs only $\mathsf{well\_scoped}$ and $\mathsf{marker\_types\_safe}$.

\emph{Progress.} If $\tj{\Gamma}{t}{\tau}$, $\mathsf{well\_scoped}\;[\,]\;t$ and $\mathsf{marker\_types\_safe}\;t$, then $t$ is a value or $t\step t'$ for some $t'$.
The proof uses a stronger open form (\texttt{progress\_open}) that, for a non-empty active marker set, additionally allows a pending $\keyword{perform}$ on an in-scope capability; at the empty set this third case is impossible, since an in-scope capability requires its marker to be in the active set.

\emph{Preservation.} If $\tj{\Gamma}{t}{\tau}$ and $t\step t'$, then $\tj{\Gamma}{t'}{\tau}$.
Ordinary redexes use the corresponding substitution-preservation lemmas; the runtime forms are introduced only by the handler rules: capabilities are typed at $\local$ data types, resumptions at $\local$ closure lifetimes, and delimiters are transparent at the public answer type $T_R$, subject to the $\noloc$ condition on their body type.

\emph{Type safety.} A well-typed program never reaches a \emph{stuck} term ($\neg\,\mathsf{value}$ and no step).
The capstone \texttt{type\_soundness} carries the invariant bundle as a premise; the bundle is proved preserved along every reduction sequence---the key lemma being preservation of the fused well-scopedness-with-closedness invariant (\texttt{step\_preserves\_ws\_rt})---and derived outright from source syntax, so no runtime-invariant assumption remains for source programs (\texttt{source\_type\_soundness}).

\emph{Determinism and execution.} One-step reduction is deterministic modulo marker renaming, and a certified step function is sound, and complete up to the fresh-marker choice (\texttt{stepf}, \texttt{Stepf.v}; \texttt{Determinism.v}), so the nondeterminism in \rname{S-Handle} is inessential.
Certified reflected deciders for value-hood, source-hood, lifetime subtyping, and the $\noloc$ premise accompany the evaluator (\texttt{Decide.v}).
This decidability of \emph{stepping} is a syntactic property, independent of the undecidability of full $F_{<:}$ \emph{subtyping} (Section~\ref{sec:discussion}).

\emph{Confinement.} \texttt{capability\_confined}: under the empty marker set a runtime capability $\cap_E^{m,\sigma}$ can never occupy a pure (delimiter-free for $m$) position---it must sit under a delimiter $\hdlr_m$ carrying its marker; multi-shot resumption can duplicate same-marker delimiters, so ``some delimiter with the same marker'' is exactly the provable form.
Transported along the dynamics, a capability reachable from a well-typed program always sits under a same-marker delimiter: the transported form (\texttt{capability\_never\_exposed}) takes $\mathsf{well\_scoped}$ along the trace as a premise, and for source programs that premise is discharged by the preserved invariant bundle (\texttt{source\_capability\_never\_exposed}), leaving only the initial typing.
Confinement is moreover \emph{structural}, not only active-position: every syntactic occurrence of a capability in every reachable state---under $\lambda$s, in constructor arguments, in stored operation bodies---has a delimiter carrying its marker in scope along its occurrence path (\texttt{source\_capability\_occurrence\_delimited}).
The corollary an algebraic-effects reader looks for first also falls out: along any execution of a well-typed source program, an active $\keyword{perform}$ is never undelimited---a delimiter carrying its capability's marker always encloses it, so the operation is handled (\texttt{source\_effect\_safety}).
Confinement is therefore a proved theorem, not an assumption.

\emph{Guarded-channel safety.} The handler boundary has two \emph{guarded channels}: the operation argument entering at a $\keyword{perform}$, and the ordinary delimited-body answer leaving through $\hdlr_m$.
The headline result is stated on \emph{executed events}: \texttt{source\_boundary\_step\_noloc} shows that every boundary transition on either channel carries a value typed at \emph{that channel's declared type}, which the typing rules force to be escapable (conclusion \texttt{boundary\_channel\_typed}).
The conclusion is tied to the fired transition itself---both endpoints of the step and its side conditions---so the typing claim cannot drift to a different decomposition of the state.
A state-shaped form (\texttt{source\_handler\_boundary\_noloc}) and a data-type specialization (\texttt{source\_boundary\_value\_non\_local}: a constructor with non-local top-level annotation) accompany it.
Three flows are deliberately unguarded.
The resumption enters the operation body with $\local$ closure lifetime (\texttt{source\_boundary\_resumption\_local}): it cannot pass either guarded channel, though it may be retained through $\local$ channels---applying it remains safe because its body re-installs the delimiter.
The operation \emph{result} delivered by the resumption may be $\local$ (Section~\ref{sec:discussion}): it lands back \emph{inside} the delimited scope, where the same guarded channels govern any further escape.
And an operation clause that never resumes consumes the delimiter---\rname{H-Perform} re-installs it only inside the resumption---so its answer, possibly $\local$ since \rname{T-Handle} intentionally leaves $T_R$ unconstrained, becomes the $\keyword{handle}$ expression's own result in the enclosing scope; in the extreme it may even hand out the resumption upcast to $\Any\,\local$, harmlessly, since $\Any$ has no eliminator and $\Any\,\local$ is not escapable---precisely why the depth theorem below is stated at \emph{escapable} types.
The two guarded channels are thus safe for local values; the boundary as a whole is deliberately \emph{not} impermeable---the unguarded flows are part of the design.

\emph{Non-escape.} The headline form is depth-blind and applies at \emph{any} escapable type: \texttt{source\_noloc\_result\_no\_runtime\_forms} shows that a source program typed at a $\noloc$ type only ever delivers values containing no capability or delimiter at any depth---the \emph{runtime} forms, not source $\keyword{handle}$ expressions, which may well sit under a $\lambda$ in a delivered closure.
Specialized to data types, \texttt{free\_data\_result\_top\_lifetime} shows that a computation staying typed at $K\,\free\,\overline{\tau}$ along every reduction terminates only in constructors whose top-level lifetime annotation is provably non-local; by preservation the trace premise disappears for source programs (\texttt{source\_free\_data\_result\_top\_lifetime}).
A subtyping-level companion, \texttt{local\_data\_not\_escapes}, shows the coercion $\tsub{\Gamma}{K\,\local\,\overline{\tau}}{K\,\free\,\overline{\tau}}$ underivable for any \emph{data} tag $K$---effect tags and $\Any$ are excluded by the statement's premises---in any runtime environment: local data cannot even be retyped as free, let alone delivered.
The type information therefore forbids the observable leakage of local values.

\emph{Guarantee suite.} The principal source-facing guarantees are bundled as one record and one theorem with a single \texttt{Print Assumptions} target (\texttt{source\_safety\_suite}, \texttt{Guarantees.v}):
\[
\mathsf{eval\_ctx}\;\Gamma \,\wedge\, \mathsf{sourceb}(t)=\mathsf{true} \,\wedge\, \tj{\Gamma}{t}{T} \;\Rightarrow\; \mathsf{source\_guarantees}\;\Gamma\;t\;T,
\]
where $\mathsf{sourceb}$ is the certified source-hood decider and the record's seven fields map the results of this section to the artifact: \texttt{sg\_type\_safety} (no reachable state is stuck), \texttt{sg\_invariants\_preserved} (the invariant bundle holds in every reachable state), \texttt{sg\_capabilities\_confined} and \texttt{sg\_capability\_occurrences\_delimited} (both confinement forms), \texttt{sg\_noloc\_results\_cap\_free} (depth non-escape of runtime forms at escapable result types), \texttt{sg\_boundary\_channels\_noloc} (every executed boundary event typed at its channel's declared, escapable type), and \texttt{sg\_boundary\_resumptions\_local} (the resumption enters its clause as a $\local$, delimiter-re-installing $\lambda$).

\subsection{Trusted base and proof engineering}

The substitution metatheory---shift and substitution commutation laws, weakening, narrowing, inversions---and \emph{all} per-redex substitution-preservation lemmas (for $\beta$, type-, and lifetime-application, \texttt{match}, and \texttt{perform}) are proved outright, as is the preservation of every runtime invariant.
The built development contains no \texttt{Axiom}, no \texttt{Admitted}, and no \texttt{Conjecture}: \texttt{Print Assumptions} on every capstone theorem of this section---including \texttt{type\_soundness}, the source-program corollaries, and \texttt{source\_safety\_suite}---reports \emph{closed under the global context}, so the trusted base is the Rocq kernel and its standard library.

The design of the well-scopedness invariant is delicate: the two natural formulations are not preserved by reduction.
A direct marker-confinement invariant breaks at the $\keyword{perform}$ contraction, which deposits stored operation clauses \emph{outside} the consumed delimiter; and marker \emph{freshness} breaks under $\beta$-substitution, which duplicates subterms and destroys uniqueness.
The predicate that survives records \emph{provenance} instead: each capability's stored operation clauses must be well-scoped at the scope outside that capability's own delimiter---exactly what the $\keyword{perform}$ contraction needs---and it is monotone under scope extension, so installing a new delimiter or moving a term under binders cannot invalidate it.
The reduction semantics still mints globally fresh markers in \rname{S-Handle}; freshness is simply not part of the \emph{preserved} invariant, and marker shadowing is harmless because the semantics always addresses the \emph{innermost} occurrence.
The development is about $2.9\times10^{4}$ lines with about $1.1\times10^{3}$ \texttt{Qed} obligations; the substitution metatheory (\texttt{meta/}) alone accounts for roughly $47\%$ of the lines, the price of pure de Bruijn indices with hand-rolled shift and substitution functions.

\section{Evaluation}\label{sec:evaluation}

The evaluation has four parts---reproducibility, worked examples, negative witnesses, and complete-term rejections---followed by a comparison with related systems and a note on practical use.

\emph{Reproducibility.} The verified core builds with standard Rocq/Coq and the Stdlib only; \texttt{make} generates a makefile from \texttt{src/\_CoqProject} and compiles the modules of Table~\ref{tab:modules}.
The development targets Rocq~9.1; exactly the files listed in \texttt{src/\_CoqProject} constitute the verified development; the example files are checked and gated, but no metatheory result depends on them.
A separate target, \texttt{make check-assumptions}, re-runs the \texttt{Print Assumptions} gate over a capstone list \emph{derived} from the sources (including the example witnesses), and \texttt{make verify} adds a documentation-freshness gate; both run in continuous integration, and the maintenance tooling additionally requires Python~3, while the verified development itself needs only Rocq.
The repository ships under the MIT license with a generated theorem index and build statistics (\texttt{STATS.md}), and the artifact guide documents the tag-release-archive procedure (GitHub release plus Zenodo DOI) to be run on the reviewed commit.

\emph{Worked examples.} \texttt{Examples.v}/\texttt{ExamplesProofs.v} mechanize the program of Listing~\ref{lst:state} in monomorphic, inlined form (\texttt{state\_example}; the listing's summing closing call is \texttt{state\_sum\_example}, machine-checked to reduce to $5$), together with reader, exception, and optionality handlers, each with a machine-checked typing proof and a reduction trace.
On the two-operation \texttt{withState} run, the capstone guarantees are witnessed directly (\texttt{ExamplesSafety.v}): soundness on the run itself, executed and typed events on \emph{both} guarded channels (\texttt{state\_example\_boundary\_return\_event}, \texttt{state\_example\_boundary\_perform\_event}), and a concrete reachable state exhibiting the capability under a same-marker delimiter, so the confinement witnesses are not vacuous (\texttt{state\_example\_cap\_reachable}).
The rejected \texttt{leak} of Listing~\ref{lst:state} fails exactly the $\noloc_\Gamma(T_B)$ premise of \rname{T-Handle}---a machine-checked rejection of the complete program (\texttt{leak\_state\_rejected}), with a positive companion typing the capability \emph{inside} the handler (\texttt{state\_capability\_typable\_inside}).
On the one-operation reader analogue, the boundary guarantee is instantiated concretely: \texttt{reader\_example\_boundary\_noloc} specializes \texttt{source\_handler\_boundary\_noloc} to its execution.
A delegating handler, whose operation clause performs on an outer capability, carries its own safety witnesses, alongside bounded sums and multi-$\keyword{perform}$ runs.
The library also exercises the resumption discipline beyond one-shot use: a multi-shot handler whose operation body resumes twice and sums both results (\texttt{multishot\_example}, with a machine-checked reduction trace evaluating to $5$) and a forwarding example that nests handlers.
No linearity restriction is imposed or needed---a resumption is an ordinary $\local$ function, so multi-shot use is ordinary application, covered by preservation.

\emph{Negative witnesses.} Three locality violations are rejected by \emph{machine-checked} computations: storing a $\local$ reader inside a $\free$ \texttt{Option} (\texttt{rejected\_testWithState}), exposing \texttt{Option}\,$\local$ through a $\free$ interface (\texttt{rejected\_crashBox}), and trapping an existential lifetime in an invariant constructor argument (\texttt{rejected\_crashEndo}).
The first two are stated on the underivability of the actual typing premise---$\noloc_\Gamma(\tau)$ for the offending interface type $\tau$---and discharged by a certified reflected decider for that judgment (\texttt{nolocb}, \texttt{Decide.v}) computing $\mathsf{false}$ under \texttt{vm\_compute}; the third computes $\elimp$ returning \emph{none} by \texttt{reflexivity}.
The first witness merits a closer look.
The offending program stores a reader capability in an \texttt{Option} and hands it out at the $\free$-annotated interface type $\mathtt{Option}\,\free\,(\mathit{Reader}\,\local\,\mathsf{Unit})$.
The lifetime restriction $\ltof$ of that type is the minimum over its lifetime annotations: the $\local$ of the stored capability type drags the whole minimum down to $\local$, and the judgment $\lsub{\Gamma}{\local}{\free}$ that $\noloc$ would then need is exactly the underivable statement behind \texttt{local\_data\_not\_escapes} (Section~\ref{sec:mechanization}).
The rejection is thus a negative verdict of a certified decider---a \texttt{Qed}-checked fact, not a hand-waved counterexample.
These computations show the checks firing on the canonical leakage channels; the following theorems lift rejection to complete programs.

\emph{Complete-term rejections and precision.} Beyond side-condition computations, the library proves rejection of \emph{complete} offending programs, each paired with a positive companion showing the underlying construction is meaningful (\texttt{ExamplesRejection.v}).
Directly returning a capability at the $\free$-annotated interface $\mathtt{Option}\,\free\,(\mathit{Reader}\,\local\,\mathsf{Unit})$ is not typable at that interface (\texttt{leak\_reader\_rejected}), while the same value is typable \emph{inside} the handler at that very type (\texttt{some\_capability\_typable\_inside}): the analysis rejects the escape, not the value.
A closure capturing a capability types at $\local$ closure lifetime and is rejected at $\free$ (\texttt{ask\_closure\_typable\_local}, \texttt{ask\_closure\_rejected\_at\_free}); nesting the capability under two constructors changes nothing (\texttt{leak\_reader2\_rejected}); and the deliberately imprecise $\Any\,\local$ interface accepts what every proper data interface rejects.
One false positive is exhibited on purpose: a closed polymorphic identity is well typed (\texttt{typed\_poly\_id}) yet conservatively fails $\noloc$, since $\ltof(\forall\ldots)=\local$.
Rejection at typing time is the static complement of the dynamic guarantees of Section~\ref{sec:mechanization}---programs that would leak do not typecheck, and programs that typecheck provably do not leak.

Table~\ref{tab:comparison} situates the work; every row refers to the cited formal system rather than the full language it comes from, and the Mech.\ column flattens very different depths---RustBelt is a far deeper semantic model than our syntactic result.
Unlike Rust and OCaml \texttt{local}, the model has first-class capability values and resumptions; unlike Effekt and capture tracking, it hides field lifetimes existentially and adds no effect-specific type former.
In theorem strength, the mechanized capture-set systems prove type soundness for their disciplines; the guarantees here are \emph{dynamic} confinement statements---occurrence-level delimitation and per-event channel typing---obtained without extending the type language.

\begin{table}[h]
\centering
\caption{Comparison with related systems (Mech.\ = mechanized metatheory; EA = type-level escape control for values; Caps = first-class handler capabilities and resumptions; Exist.\ = existential field lifetimes; Former = adds an effect-specific type former such as rows, capture sets, or boxes; ``---'' = no effect system in the cited formalism).}
\label{tab:comparison}
\footnotesize
\begin{tabular}{@{}p{0.30\linewidth}ccccc@{}}
\toprule
\textbf{System} & \textbf{Mech.} & \textbf{EA} & \textbf{Caps} & \textbf{Exist.} & \textbf{Former} \\
\midrule
Rust/RustBelt~\cite{matsakis2014rust,jung2018rustbelt} & partial & yes & no & no & --- \\
OCaml \texttt{local}~\cite{lorenzen2024oxidizing} & no & yes & no & no & --- \\
Effekt~\cite{brachthauser2020effects,brachthauser2022effects} & yes & yes & boxed & no & yes \\
Capture tracking~\cite{boruch2023capturing} & yes & yes & no & no & yes \\
$Core_\Delta$~\cite{stoyan2025existential} & no & yes & no & yes & no \\
This work & yes & yes & yes & yes & no \\
\bottomrule
\end{tabular}
\end{table}

A practical use is the design of a compiler intermediate language where stack or region allocation must coexist with capability-safe handlers: the typing rules give local checks at function boundaries, the semantics fixes the runtime invariant a correct capture/resume implementation must maintain, and---with no memory model---the artifact is a candidate \emph{typing} specification, the memory or stack argument remaining separate.
Lexically scoped handlers have been compiled efficiently to regions~\cite{mueller2023regions}, which reinforces this reading.

\section{Limitations and Future Work}\label{sec:discussion}

The metatheory is closed for the calculus as stated---no axioms, admitted obligations, or unproved metatheoretic assumptions remain---so the limitations are those of the calculus itself.
First, operation typing is asymmetric: the argument and the perform-site type instantiations must be non-local, while the result may be local (so resumptions can deliver local values back to the operation site); lifting the remaining argument restriction would require threading the argument lifetime through the continuation type.
Confinement is per-marker, not per-operation: all operations of one handler share its marker and its public answer type $T_R$.

Second, there is no memory model: connecting the lexical discipline to actual stack or region allocation, and widening the two-point lattice to multi-level regions, are the main semantic extensions.
A stack-based implementation additionally requires a one-shot restriction or frame copying for resumptions---a continuation-use concern orthogonal to the value discipline proved here~\cite{tang2024soundly}.
The lattice machinery is already largely parametric: lifetime subtyping, $\elimp$, \rname{T-Match}, and the substitution tier work over arbitrary lifetime expressions (\rname{T-Match} bounds its fresh variables by an arbitrary $\Delta$); only the constant alphabet and the $\free$ target of $\noloc$ are two-point-specific.
The analysis is also deliberately coarse: $\ltof$ classifies every quantified type as $\local$, a sound but conservative choice that rejects some safe polymorphic values.

Finally, the typing relation is specified but not yet realized as a verified algorithm.
The escape-specific checks are already executable---$\elimp$ is a computable partial function, and certified reflected deciders for lifetime subtyping and the $\noloc$ premise ship with the artifact (\texttt{lt\_subb}, \texttt{nolocb}; \texttt{Decide.v})---but subtyping features full $F_{<:}$ bounded quantification with declarative transitivity, for which subtyping is undecidable~\cite{pierce1994bounded}, so a complete terminating checker is impossible in general: restricting quantifier bounds to the kernel discipline, or extracting a checker that is sound but incomplete for subtyping, are the concrete next steps.

\section{Conclusion}\label{sec:conclusion}

We presented a machine-checked metatheory of $Core_\Delta$ extended with multi-operation algebraic effect handlers, mechanizing the full system---syntax, lifetime and type subtyping, typing with a constructor minimum rule and variance-aware elimination of existential lifetimes, and a marker-indexed reduction semantics---and proving progress, preservation, and type safety with no runtime-invariant assumptions for source programs, together with occurrence-level capability confinement, guarded-channel safety, and depth non-escape of runtime capabilities and delimiters.
This discharges requirements (i)--(iv) of Section~\ref{sec:related-work}.
An executable example library and machine-checked negative witnesses show the analysis on reader/state/exception/optionality handlers and on programs it must reject.
The artifact is useful as a checked typing specification---with executable reduction semantics and escape-side-condition deciders---for compiler intermediate languages with first-class capability values and capability-safe effect handlers; widening the lifetime lattice to multi-level regions, connecting to a memory model, and extracting a certified checker are the natural next steps.

% Journal house style (see tmp/spiiras_latex/Template.tex): hand-formatted
% numbered lists in both languages, generated from bib.bib in citation order
% by gen_refs.py. Regenerate after changing citations; do not edit by hand.
% AUTO-GENERATED from bib.bib by gen_refs.py (citation order); edit the .bib, not this file.
\begin{thebibliographyEng}{99}

\bibitem{park1992escape}
Park Y.G., Goldberg B. Escape analysis on lists. Proceedings of the ACM SIGPLAN 1992 conference on Programming language design and implementation. 1992. pp. 116--127.

\bibitem{blanchet1999escape}
Blanchet B. Escape analysis for object-oriented languages: application to Java. \emph{Acm Sigplan Notices.} 1999. vol. 34. no. 10. pp. 20--34.

\bibitem{matsakis2014rust}
Matsakis N.D., Klock F.S. The rust language. Proceedings of the 2014 ACM SIGAda annual conference on High integrity language technology. 2014. pp. 103--104.

\bibitem{lorenzen2024oxidizing}
Lorenzen A., White L., Dolan S., Eisenberg R.A., Lindley S. Oxidizing OCaml with modal memory management. \emph{Proceedings of the ACM on Programming Languages.} 2024. vol. 8. no. ICFP. pp. 485--514.

\bibitem{georges2025data}
Georges A.L., Peters B., Elbeheiry L., White L., Dolan S., Eisenberg R.A., Casinghino C., Pottier F., Dreyer D. Data Race Freedom {\`a} la Mode. \emph{Proceedings of the ACM on Programming Languages.} 2025. vol. 9. no. POPL.

\bibitem{akhin2021kotlin}
Akhin M., Belyaev M. Kotlin language specification. \emph{Kotlin Language Specification.} 2021.

\bibitem{jung2018rustbelt}
Jung R., Jourdan J., Krebbers R., Dreyer D. RustBelt: Securing the Foundations of the Rust Programming Language. \emph{Proceedings of the ACM on Programming Languages.} 2018. vol. 2. no. POPL. pp. 1--34.

\bibitem{jung2020stacked}
Jung R., Dang H., Kang J., Dreyer D. Stacked Borrows: An Aliasing Model for Rust. \emph{Proceedings of the ACM on Programming Languages.} 2020. vol. 4. no. POPL. pp. 1--32.

\bibitem{astrauskas2022prusti}
Astrauskas V., B{\'{i}}l{\'{y}} A., Fiala J., Grannan Z., Matheja C., M{\"{u}}ller P., Poli F., Summers A.J. The Prusti Project: Formal Verification for Rust. NASA Formal Methods - 14th International Symposium, NFM 2022. 2022. pp. 88--108.

\bibitem{denis2022creusot}
Denis X., Jourdan J., March{\'{e}} C. Creusot: A Foundry for the Deductive Verification of Rust Programs. Formal Methods and Software Engineering - 23rd International Conference on Formal Engineering Methods, ICFEM 2022. 2022. pp. 90--105.

\bibitem{matsushita2022rusthornbelt}
Matsushita Y., Denis X., Jourdan J., Dreyer D. RustHornBelt: a semantic foundation for functional verification of Rust programs with unsafe code. PLDI '22: 43rd ACM SIGPLAN International Conference on Programming Language Design and Implementation. 2022. pp. 841--856.

\bibitem{gaher2024refinedrust}
G{\"{a}}her L., Sammler M., Jung R., Krebbers R., Dreyer D. RefinedRust: A Type System for High-Assurance Verification of Rust Programs. \emph{Proceedings of the ACM on Programming Languages.} 2024. vol. 8. no. PLDI. pp. 1115--1139.

\bibitem{plotkin2013handling}
Plotkin G.D., Pretnar M. Handling algebraic effects. \emph{Logical methods in computer science.} 2013. vol. 9.

\bibitem{brachthauser2020effects}
Brachth{\"a}user J.I., Schuster P., Ostermann K. Effects as capabilities: effect handlers and lightweight effect polymorphism. \emph{Proceedings of the ACM on Programming Languages.} 2020. vol. 4. no. OOPSLA. pp. 1--30.

\bibitem{brachthauser2022effects}
Brachth{\"a}user J.I., Schuster P., Lee E., Boruch-Gruszecki A. Effects, capabilities, and boxes: from scope-based reasoning to type-based reasoning and back. \emph{Proceedings of the ACM on Programming Languages.} 2022. vol. 6. no. OOPSLA1. pp. 1--30.

\bibitem{boruch2023capturing}
Boruch-Gruszecki A., Odersky M., Lee E., Lhot{\'a}k O., Brachth{\"a}user J. Capturing types. \emph{ACM Transactions on Programming Languages and Systems.} 2023. vol. 45. no. 4. pp. 1--52.

\bibitem{xie2021generalized}
Xie N., Leijen D. Generalized evidence passing for effect handlers: efficient compilation of effect handlers to C. \emph{Proceedings of the ACM on Programming Languages.} 2021. vol. 5. no. ICFP. pp. 1--30.

\bibitem{xie2022first}
Xie N., Cong Y., Ikemori K., Leijen D. First-class names for effect handlers. \emph{Proceedings of the ACM on Programming Languages.} 2022. vol. 6. no. OOPSLA2. pp. 30--59.

\bibitem{stoyan2025existential}
Stoyan A.S. Type-based Escape Analysis with Existential Lifetimes. \emph{Trudy ISP RAN / Proc. ISP RAS.} 2026. vol. 38. no. 1. pp. 71--78.

\bibitem{tofte1997region}
Tofte M., Talpin J. Region-Based Memory Management. \emph{Information and Computation.} 1997. vol. 132. no. 2. pp. 109--176.

\bibitem{fluet2006monadic}
Fluet M., Morrisett G. Monadic Regions. \emph{Journal of Functional Programming.} 2006. vol. 16. no. 4--5. pp. 485--545.

\bibitem{zyuzin2021contextual}
Zyuzin N., Nanevski A. Contextual modal types for algebraic effects and handlers. \emph{Proceedings of the ACM on Programming Languages.} 2021. vol. 5. no. ICFP. pp. 1--29.

\bibitem{van2022handling}
van der Rest C., Reinders J., Poulsen C.B. Handling Higher-Order Effects. \emph{arXiv preprint arXiv:2203.03288.} 2022.

\bibitem{van2024framework}
van den Berg B., Schrijvers T. A framework for higher-order effects \& handlers. \emph{Science of Computer Programming.} 2024. vol. 234. pp. 103086.

\bibitem{tang2025modal}
Tang W., White L., Dolan S., Hillerstr{\"o}m D., Lindley S., Lorenzen A. Modal effect types. \emph{Proceedings of the ACM on Programming Languages.} 2025. vol. 9. no. OOPSLA1. pp. 1130--1157.

\bibitem{sivaramakrishnan2021retrofitting}
Sivaramakrishnan K.C., Dolan S., White L., Kelly T., Jaffer S., Madhavapeddy A. Retrofitting effect handlers onto OCaml. PLDI '21: 42nd ACM SIGPLAN International Conference on Programming Language Design and Implementation. 2021. pp. 206--221.

\bibitem{tang2024soundly}
Tang W., Hillerstr{\"o}m D., Lindley S., Morris J.G. Soundly Handling Linearity. \emph{Proceedings of the ACM on Programming Languages.} 2024. vol. 8. no. POPL.

\bibitem{vanrooij2025affect}
van Rooij O., Krebbers R. Affect: An Affine Type and Effect System. \emph{Proceedings of the ACM on Programming Languages.} 2025. vol. 9. no. POPL.

\bibitem{wu2014effect}
Wu N., Schrijvers T., Hinze R. Effect handlers in scope. Proceedings of the 2014 ACM SIGPLAN Symposium on Haskell. 2014. pp. 1--12.

\bibitem{biernacki2019binders}
Biernacki D., Pir{\'o}g M., Polesiuk P., Sieczkowski F. Binders by day, labels by night: effect instances via lexically scoped handlers. \emph{Proceedings of the ACM on Programming Languages.} 2019. vol. 4. no. POPL. pp. 1--29.

\bibitem{odersky2021safer}
Odersky M., Boruch-Gruszecki A., Brachth{\"a}user J.I., Lee E., Lhot{\'a}k O. Safer exceptions for Scala. Proceedings of the 12th ACM SIGPLAN International Symposium on Scala. 2021. pp. 1--11.

\bibitem{odersky2022scoped}
Odersky M., Boruch-Gruszecki A., Lee E., Brachth{\"a}user J., Lhot{\'a}k O. Scoped capabilities for polymorphic effects. \emph{arXiv preprint arXiv:2207.03402.} 2022.

\bibitem{KEEP-context-parameters}
Serrano A., Akhin M., Bobko N., Gorbunov I., Zarechenskii M., Zharkov D. Context parameters (KEEP-367). 2024. Available at: \url{}.

\bibitem{osvald2016gentrification}
Osvald L., Essertel G., Wu X., Alay{\'o}n L.I.G., Rompf T. Gentrification gone too far? affordable 2nd-class values for fun and (co-) effect. \emph{ACM SIGPLAN Notices.} 2016. vol. 51. no. 10. pp. 234--251.

\bibitem{leijen2017type}
Leijen D. Type directed compilation of row-typed algebraic effects. Proceedings of the 44th ACM SIGPLAN Symposium on Principles of Programming Languages. 2017. pp. 486--499.

\bibitem{zhang2016accepting}
Zhang Y., Salvaneschi G., Beightol Q., Liskov B., Myers A.C. Accepting blame for safe tunneled exceptions. \emph{ACM SIGPLAN Notices.} 2016. vol. 51. no. 6. pp. 281--295.

\bibitem{zhang2020handling}
Zhang Y., Salvaneschi G., Myers A.C. Handling bidirectional control flow. \emph{Proceedings of the ACM on Programming Languages.} 2020. vol. 4. no. OOPSLA. pp. 1--30.

\bibitem{xu2023formalizing}
Xu Y., Odersky M. Formalizing Box Inference for Capture Calculus. \emph{arXiv preprint arXiv:2306.06496.} 2023.

\bibitem{xu2025s}
Xu Y., Bra{\v{c}}evac O., Pham C.N., Odersky M. What’s in the Box?. \emph{Ergonomic and Expressive Capture Tracking over Generic Data Structures (Extended Version),(Aug. 2025). doi.} 2025. vol. 10. pp. 17.

\bibitem{stoyan2024modern}
Stoyan A. Modern Effect Systems Overview. 2024 Ivannikov Ispras Open Conference (ISPRAS). 2024. pp. 1--7.

\bibitem{aydemir2005poplmark}
Aydemir B.E., Bohannon A., Fairbairn M., Foster J.N., Pierce B.C., Sewell P., Vytiniotis D., Washburn G., Weirich S., Zdancewic S. Mechanized Metatheory for the Masses: The POPLMark Challenge. Theorem Proving in Higher Order Logics. 2005. pp. 50--65.

\bibitem{aydemir2008engineering}
Aydemir B.E., Chargu{\'e}raud A., Pierce B.C., Pollack R., Weirich S. Engineering Formal Metatheory. \emph{ACM SIGPLAN Notices.} 2008. vol. 43. no. 1. pp. 3--15.

\bibitem{sewell2010ott}
Sewell P., Nardelli F.Z., Owens S., Peskine G., Ridge T., Sarkar S., Strni{\v{s}}a R. Ott: Effective Tool Support for the Working Semanticist. \emph{Journal of Functional Programming.} 2010. vol. 20. no. 1. pp. 71--122.

\bibitem{schaefer2015autosubst}
Sch{\"a}fer S., Tebbi T., Smolka G. Autosubst: Reasoning with de Bruijn Terms and Parallel Substitutions. Interactive Theorem Proving. 2015. pp. 359--374.

\bibitem{stark2019autosubst2}
Stark K., Sch{\"a}fer S., Kaiser J. Autosubst 2: Reasoning with Multi-sorted de Bruijn Terms and Vector Substitutions. Proceedings of the 8th ACM SIGPLAN International Conference on Certified Programs and Proofs. 2019. pp. 166--180.

\end{thebibliographyEng}


\begin{aboutAuthors}
\textbf{Andrey Stoyan}~--- researcher in type systems; National Research University Higher School of Economics (HSE University), Saint Petersburg, Russian Federation.
Position: PhD student.
ORCID: 0009-0004-6935-1447.
E-mail: a.stoyan@hse.ru.
Research interests: type systems and effect systems.\smallskip

\textbf{Acknowledgements.} This research received no external funding.
The Claude Opus 4.8 and Claude Fable 5 models (Anthropic) were used as assistants during the development of the artifact and the preparation of the manuscript; all proofs are machine-checked and all content was reviewed by the author.
\end{aboutAuthors}

\newpage
\selectlanguage{russian}
\makerustitle
\normalsize

\begin{aboutAuthors}
\textbf{Стоян Андрей}~--- исследователь в области систем типов; Национальный исследовательский университет <<Высшая школа экономики>> (НИУ ВШЭ), Санкт-Петербург, Российская Федерация.
Должность: аспирант.
ORCID: 0009-0004-6935-1447.
E-mail: a.stoyan@hse.ru.
Область научных интересов: системы типов, системы эффектов.\smallskip

\textbf{Поддержка исследований.} Финансирование исследования из внешних источников не привлекалось.
При разработке артефакта и подготовке рукописи использовались модели Claude Opus 4.8 и Claude Fable 5 (Anthropic); все доказательства проверены машиной, весь текст выверен автором.
\end{aboutAuthors}

\input{escape-analysis-proofs-refs-rus}

\end{document}
